\documentclass[11pt]{amsart}
\usepackage[T1]{fontenc}
\usepackage[utf8]{inputenc}
\usepackage{lmodern}
\usepackage{amsmath,amssymb,amsthm,mathtools}
\usepackage{enumitem}
\usepackage[margin=1in]{geometry}
\usepackage{microtype}
\usepackage[hidelinks]{hyperref}
\usepackage{cite}
\usepackage{xcolor}

\newtheorem{theorem}{Theorem}[section]
\newtheorem{lemma}[theorem]{Lemma}
\newtheorem{proposition}[theorem]{Proposition}
\newtheorem{corollary}[theorem]{Corollary}
\newtheorem{conjecture}[theorem]{Conjecture}
\theoremstyle{definition}
\newtheorem{definition}[theorem]{Definition}
\newtheorem{remark}[theorem]{Remark}
\newtheorem{example}[theorem]{Example}

\newcommand{\F}{\mathbb F}
\newcommand{\N}{\mathbb N}
\newcommand{\PP}{\mathbb P}
\newcommand{\Aff}{\mathbb A}
\newcommand{\ev}{\operatorname{ev}}
\newcommand{\PRM}{\operatorname{PRM}}
\newcommand{\RM}{\operatorname{RM}}
\newcommand{\Span}{\operatorname{Span}}
\newcommand{\supp}{\operatorname{supp}}
\newcommand{\Sym}{\operatorname{Sym}}
\newcommand{\gr}{\operatorname{gr}}
\newcommand{\ini}{\operatorname{in}}
\newcommand{\rank}{\operatorname{rank}}
\newcommand{\card}[1]{\left|#1\right|}
\newcommand{\sq}[1]{\left\lfloor q^{#1}\right\rfloor}
\numberwithin{equation}{section}
\setlength{\emergencystretch}{2em}

\newcommand{\rojo}{\textcolor{red}}

\title[Common zeros of homogeneous polynomials]
{Common zeros of homogeneous polynomials and the weight hierarchy of projective Reed--Muller codes}
\author{Rodrigo San-Jos\'e}
\address{Department of Mathematics, Virginia Tech, Blacksburg, VA, USA}
\email{rsanjose@vt.edu}
\date{}
\thanks{The author was partially supported by the NSF grant DMS-2401558, the Commonwealth Cyber Initiative, an AMS-Simons Travel Grant, and by Grant PID2022-138906NB-C21 funded by MICIU/AEI/10.13039/501100011033 and by ERDF/EU}
\subjclass[2020]{14G15, 13D40, 94B27, 11T71}
\keywords{Rational points, Hilbert functions, Projective Reed--Muller codes, generalized Hamming weights, footprint bounds}

\begin{document}
\begin{abstract}
We determine the maximum number of common zeros in
projective space of any set of linearly independent
homogeneous polynomials of the same degree, when the degree is
smaller than the cardinality of the finite field. This proves the
conjecture of Beelen, Datta and Ghorpade and determines the complete
weight hierarchy of the corresponding projective Reed--Muller
codes. The proof uses Green's hyperplane restriction bound,
combined with a footprint argument that reduces the problem
to the affine case solved by Heijnen and Pellikaan.
Wei duality then determines the hierarchy
in sufficiently large degrees, recovering the complete hierarchy
in every degree for the projective plane.
Finally, we prove an upper bound valid in every degree, and conjecture that it
is always sharp.
\end{abstract}
\maketitle

\section{Introduction}\label{sec:intro}
Let $q$ be a prime power and let $m,d\ge1$. We write
$S=\F_q[X_0,\ldots,X_m]$ and denote by $S_d$ the vector space
of homogeneous polynomials of degree $d$ in $S$. For
$F_1,\ldots,F_r\in S_d$, let $V(F_1,\ldots,F_r)$ be their
common zero set in $\PP^m$, and let
$V(F_1,\ldots,F_r)(\F_q)$ denote its set of $\F_q$-rational
points. We are interested in the maximum number of such points
when the polynomials are linearly independent. More precisely, for
$1\le r\le\binom{m+d}{d}$, we set
\[
 e_r(d,m):=\max\big\{\card{V(F_1,\ldots,F_r)(\F_q)}:
 F_1,\ldots,F_r\in S_d\text{ are linearly independent}\big\}.
\]
We write $p_j=1+q+\cdots+q^j$ for $j\ge0$, and $p_j=0$ for $j<0$.
The problem of interest for this paper is to determine $e_r(d,m)$ in terms of
$q,d,m$ and $r$.

This problem is closely connected with projective Reed--Muller
codes, introduced by Lachaud~\cite{Lachaud}. The code
$\PRM_q(d,m)$ is obtained by evaluating the elements of $S_d$
at fixed representatives of the points of $\PP^m(\F_q)$. They have received considerable attention recently \cite{sanjose_decoding_PRM,sanjoseHullsPRM,sanjoseHullvariationPRM,sanjoseSSCPRM,KaplanKim2025,SongYue2026,SongLuo2026}.
Its basic parameters were studied by Lachaud~\cite{Lachaud1990}
and S{\o}rensen~\cite{Sorensen}; see also~\cite{GL24}.
Generalized Hamming weights, studied systematically by
Wei~\cite{Wei}, extend the minimum distance by considering the
smallest support of a subcode of each possible dimension.
For $1\le d<q$, the evaluation map on $S_d$ is injective, and
the $r$-th generalized Hamming weight satisfies
\begin{equation}\label{eq:weightszeros}
 d_r(\PRM_q(d,m))=p_m-e_r(d,m).
\end{equation}
Thus determining $e_r(d,m)$ for all $r$ in this range is
equivalent to determining the complete weight hierarchy of
$\PRM_q(d,m)$.

Computing the weight hierarchy of an arbitrary linear code is
NP-hard, since this includes computing its minimum
distance~\cite{vardyIntractability}.
Exact values and bounds have been obtained for many families
of codes with additional structure ~\cite{munueraGHWhermitica,sanjoseGHWNT,sanjose_complete_intersections,sanjoseGHWMPC,beelenGHWcartesian,sanjose_squarefree}. For projective Reed--Muller codes, the case of one equation goes back to a conjecture of Tsfasman,
proved independently by Serre~\cite{Serre} and
S{\o}rensen~\cite{Sorensen} in 1991. Their result gives
\[
 e_1(d,m)=dq^{m-1}+p_{m-2}\qquad\text{for }1\le d\le q.
\]
Boguslavsky~\cite{Boguslavsky1997} determined $e_2(d,m)$ for
$d<q-1$, and Zanella~\cite{Zanella} settled the case of quadratic
equations for every admissible $r$. The affine version of the
problem was solved by Heijnen and Pellikaan~\cite{HP}, who
determined all the generalized Hamming weights of affine
Reed--Muller codes.

A proposed formula for $e_r(d,m)$ for all $r$ and $d<q-1$
became known as the Tsfasman--Boguslavsky conjecture.
Datta and Ghorpade proved it for $r\le m+1$~\cite{DG}, but
also showed that it can fail for larger $r$~\cite{DGRemarks2017}.
They proposed a new formula for
$1\le r\le\binom{m+d-1}{d-1}$ and $d<q$
in~\cite[Conj.~6.6]{DG}. Beelen, Datta and Ghorpade established
this formula for further values of $r$~\cite{BDG18}, and then
formulated a complete conjecture covering every admissible $r$
when $d<q$~\cite[Conj.~2.4]{BDG22}.
They showed that the conjectured value is always a lower bound
and obtained upper bounds using the projective footprint
bound and extremal combinatorics~\cite{BDG19footprint,BDG22}.

More recently, Singhal and Lin~\cite{SL} proved the complete
conjecture when $q$ is sufficiently large in terms of $m,d,r$.
Lin and Singhal~\cite[Thm.~1.5]{LinSinghal2025} settled the
plane case $m=2$ for every $d<q$, through a related problem
about zero-dimensional intersections of hypersurfaces. For this related problem,
Lin~\cite[Thm.~1.5]{LinArtinian2026} proved the conjectured
formula for a range of $r$. The conjecture for these zero-dimensional intersections
also follows from the Eisenbud--Green--Harris conjecture \cite{EisenbudGreenHarris1993,eisenbud_green_harris};
see~\cite[Sec.~2.2 and Prop.~3.1]{LinDimension2026}.
A different approach, based on a recursive construction of
projective Reed--Muller codes, gives bounds for their
generalized Hamming weights~\cite{SanJose}.

In this paper we prove the complete conjecture of Beelen,
Datta and Ghorpade for every $d<q$. We state it in the
equivalent form given in~\cite[Cor.~2.6]{BDG22}, using Macaulay
expansions. Recall that every nonnegative integer $h$ has a
unique Macaulay expansion of degree $d$
\begin{equation}\label{eq:macaulay}
 h=\sum_{i=1}^d\binom{\mu_i+i}{i},
 \qquad -1\le\mu_1\le\cdots\le\mu_d.
\end{equation}
We use the convention $\binom ab=0$ if $a<b$ and $b\ge0$.
In particular, $h=0$ corresponds to $\mu_1=\cdots=\mu_d=-1$.
The equivalent form $h=\sum_{i=1}^d\binom{\nu_i}{i}$, with
$0\le\nu_1<\cdots<\nu_d$, is obtained by setting
$\nu_i=\mu_i+i$.

\begin{theorem}[Beelen--Datta--Ghorpade conjecture]\label{thm:main}
Let $q$ be a prime power, $m\ge1$, $1\le d<q$, and
$1\le r\le N=\binom{m+d}{d}$. Let $(\mu_1,\ldots,\mu_d)$ be the
Macaulay tuple of $N-r$ as in~\eqref{eq:macaulay}. Then
\begin{equation}\label{eq:mainzeros}
 e_r(d,m)=p_{\mu_d}+\sum_{i=1}^{d-1}\sq{\mu_i}.
\end{equation}
Consequently,
\begin{equation}\label{eq:mainweights}
 d_r(\PRM_q(d,m))
 =p_m-p_{\mu_d}-\sum_{i=1}^{d-1}\sq{\mu_i}.
\end{equation}
\end{theorem}
Combined with Wei duality and S{\o}rensen's description of the dual code \cite{Sorensen},
Theorem~\ref{thm:main} also determines the weight hierarchy whenever
$d>(m-1)(q-1)$; see Corollary~\ref{cor:highdegree}.
For $m=2$, these two ranges cover every positive degree. Thus we recover
the complete plane case, which already follows from the result of
Lin and Singhal~\cite[Thm.~1.5]{LinSinghal2025} and duality. We also give an upper bound for every degree, which we conjecture is always sharp. 

The main step in the proof is a comparison between projective
and affine point sets, using Hilbert functions and footprint
techniques~\cite{geilbezout,BDG19footprint}. Starting from a
set of points $Z$, we construct a zero-dimensional
quotient whose vector-space dimension is $\card Z$.
Among its standard monomials, those with all exponents less than
$q$ give a decreasing set of exponent vectors, which can be
realized by an affine rational point set.
We construct sufficiently many independent polynomials vanishing
on this affine set, and we control the difference between the
cardinalities of the two point sets. The Heijnen--Pellikaan theorem then gives the
required upper bound for $e_r(d,m)$.

To obtain this comparison, we work over an extension
field and use linear forms whose coefficients are successive
$q$th powers. We prove that Green's hyperplane restriction bound~\cite{Green,Caviglia}
remains valid for successive quotients by these linear forms
in degrees less than $q$.
Together with the equations satisfied by the rational points,
this provides additional monomials in the initial ideal and
strengthens the resulting footprint bound. 
Related refinements of footprint estimates appear in work on
norm-trace codes~\cite{sanjoseGHWNT,geilbezout} and weighted projective
spaces~\cite{sanjose_conjecture}. For the upper bound on generalized Hamming weights in arbitrary degree, we mainly use different decompositions of the rational points of $\PP^m(\F_q)$ into affine and projective spaces, similar to the arguments in \cite{SanJose,sanjose_recursive_WPRM}.

The paper is organized as follows. Section~\ref{sec:prelim} fixes
notation and recalls Wei duality, the affine common-zero bound, and Green's theorem.
In Section~\ref{sec:restriction} we prove that we can repeatedly use Green's theorem when taking quotients by the chosen linear forms. Section~\ref{sec:transfer} contains the footprint
comparison with affine decreasing sets, the proof of Theorem~\ref{thm:main}
and its consequence for high degrees.
In Section~\ref{sec:chain} we give a nested
family attaining all the weights and the resulting recursion.
Section~\ref{sec:general} gives a bound in arbitrary
degree, conjectures that it is always sharp, and determines
some generalized Hamming weights for two infinite families of
projective Reed--Muller codes.
The recursive formula used to compute this bound is proved
in Appendix~\ref{app:hilbertfunctions}.



\section{Preliminaries}\label{sec:prelim}

Throughout, $\N$ denotes the set of nonnegative integers. If $\Bbbk$ is a
field and $R$ is a standard graded $\Bbbk$-algebra, we write its Hilbert function as
$H_R(t)=\dim_{\Bbbk} R_t$. When the ambient ring $R$ is clear, we write
$(f_1,\ldots,f_s)$ for the ideal
$f_1R+\cdots+f_sR$. Similarly, $(I,f_1,\ldots,f_s)$
denotes the ideal of $R$ generated by $I$ and
$f_1,\ldots,f_s$, extending $I$ to $R$ when necessary. We also write $IR$ when we want to make the ambient ring explicit.

\subsection{Evaluation codes and Hilbert functions}

For a linear subspace $D\subseteq\F_q^n$, its support is
\[
 \supp(D):=\{i: c_i\ne0\text{ for some }c\in D\}.
\]
If $C$ is a linear code and $1\le r\le\dim C$, its $r$-th generalized
Hamming weight is
\[
 d_r(C):=\min\{\card{\supp(D)}:D\subseteq C,\ \dim D=r\}.
\]
The sequence of these numbers is the weight hierarchy of $C$.
Write $C^\perp$ for the dual code with respect to the standard inner
product on $\F_q^n$. The following theorem shows that the weight
hierarchy of a code determines that of its dual.

\begin{theorem}[Wei duality {\cite[Thm.~3]{Wei}}]\label{thm:weiduality}
Let $C\subseteq\F_q^n$ be a linear code of dimension $k$. Then
\[
 \{d_r(C):1\le r\le k\}
 =\{1,\ldots,n\}\setminus
   \{n+1-d_s(C^\perp):1\le s\le n-k\}.
\]
\end{theorem}

For each point of $\PP^m(\F_q)$, choose the representative whose leftmost nonzero coordinate is one. Evaluation at these representatives gives
\[
 \ev:S_d\longrightarrow\F_q^{p_m},
 \qquad F\longmapsto(F(P))_{P\in\PP^m(\F_q)}.
\]
Its image is the projective Reed--Muller code $\PRM_q(d,m)$. If
$d<q$, this map is injective, which proves both $\dim\PRM_q(d,m)=\binom{m+d}{d}$ and
the relation~\eqref{eq:weightszeros}.

For a finite set $Z\subseteq\PP^m(\F_q)$, let $I(Z)$ denote its
homogeneous vanishing ideal, and we denote 
$$H_Z(t):=H_{S/I(Z)}(t)=\dim_{\F_q}(S/I(Z))_t.$$

Since
\[
 \dim_{\F_q}I(Z)_d=\binom{m+d}{d}-H_Z(d),
\]
the set $Z$ is contained in the common zero set
of $r$ linearly independent elements of $S_d$ if and only if
\begin{equation}\label{eq:HZbound}
 H_Z(d)\le\binom{m+d}{d}-r.
\end{equation}
Taking the maximum over such sets gives
\[
 e_r(d,m)=
 \max\left\{
 \card Z: Z\subseteq\PP^m(\F_q),\
 H_Z(d)\le\binom{m+d}{d}-r
 \right\},
\]
which gives the connection between Hilbert functions and $e_r(d,m)$, and, consequently, with the generalized Hamming weights of projective Reed--Muller codes.
The affine Reed--Muller code $\RM_q(d,m)$ is obtained in a similar fashion by evaluating
the polynomials of total degree at most $d$ on $\F_q^m$.

\subsection{Macaulay expansions and Green's theorem}

For $h$ as in~\eqref{eq:macaulay}, let 
\[
 g_d(h):=\sum_{i=1}^d\binom{\mu_i+i-1}{i}.
\]
The function $g_d$ is nondecreasing. Its $s$-fold iterate satisfies
\[
 g_d^{\circ s}(h)=\sum_{i=1}^d\binom{\mu_i+i-s}{i}.
\]
Indeed, the nonzero terms on the right have strictly increasing upper
indices and form a Macaulay expansion. The formula follows by induction on $s$. We will use
the following consequence: 
\begin{equation}\label{eq:vanishthreshold}
 h<\binom{d+s}{d}\quad\Longrightarrow\quad
 g_d^{\circ s}(h)=0.
\end{equation}
To see this, it is enough to note that the inequality implies $\mu_d<s$.
We also recall that
\begin{equation}\label{eq:macrange}
 \binom{d+\mu_d}{d}\le h<\binom{d+\mu_d+1}{d}.
\end{equation}

We use Green's hyperplane restriction theorem in the following form.
It holds in arbitrary characteristic (see comments after \cite[Thm. 1.4]{Caviglia}). 

\begin{theorem}[Green {\cite{Green,Caviglia}}]\label{thm:green}
Let $R$ be a standard graded algebra over an infinite field $\Bbbk$,
and let $\ell\in R_1$ be a general linear form. Then
\[
 \dim_{\Bbbk}(R/\ell R)_d\le g_d(\dim_{\Bbbk} R_d).
\]
\end{theorem}

For later use, we also recall a bound on the number of common
zeros of affine polynomials of degree less than $q$.
The following formulation is obtained from~\cite[Thm. 5.10]{HP}
and~\cite[Thm. 2.1 and Lem. 2.5]{BDG22}.

\begin{theorem}\label{thm:affine}
Let $1\le d<q$, $1\le r\le\binom{m+d}{d}$, and let
$(\mu_1,\ldots,\mu_d)$ be the Macaulay tuple of
$\binom{m+d}{d}-r$.
If $Y\subseteq\Aff^m(\F_q)$ is contained in the common zero set
of $r$ linearly independent polynomials in $\F_q[X_1,\ldots,X_m]$
of total degree at most $d$, then
\[
 \card Y\le\sum_{i=1}^d\sq{\mu_i}.
\]
\end{theorem}

\section{Change of coordinates}\label{sec:restriction}
In Section~\ref{sec:transfer}, we will use a suitable change of
coordinates in our footprint argument. We choose the new
coordinates using linear forms $\ell_j=\sum_{i=0}^m a_i^{q^j}X_i$ whose coefficients are successive $q$-th powers. This choice gives polynomial relations coming
from the $\F_q$-rationality of the points, which will help us
identify monomials in the initial ideal. In particular, $\ell_0$ avoids all the rational points of $\PP^m(\F_q)$, which allows us to quotient by it and still recover the correct number of points. Note that we have $\ell_j(P)=\ell_0(P)^{q^j}$ for any $P\in \PP^m(\F_q)$.

We work over an extension field so that the coefficients
of the first linear form can be algebraically independent,
without changing Hilbert functions. Green's theorem bounds
the Hilbert function after quotienting by a general linear
form. We need the same bound at each successive quotient
by $\ell_0,\ldots,\ell_m$. However, these forms are determined
by the coefficients of $\ell_0$, so we cannot simply treat
them as successive general linear forms. The purpose of
this section is to prove that, in degrees less than $q$,
they nevertheless satisfy this bound. Iterating it will
give Corollary~\ref{cor:flagvanishing}, which we use in
Section~\ref{sec:transfer} to identify further monomials
in the initial ideal.

\begin{proposition}\label{prop:frobenius}
Let $I\subseteq\F_q[X_0,\ldots,X_m]$ be a homogeneous ideal,
where $m\ge1$, and let $1\le d<q$. Set $K=\F_q(a_0,\ldots,a_m)$, $S_K=K[X_0,\ldots,X_m]$ and $\ell_j=\sum_{i=0}^m a_i^{q^j}X_i$ for $0\le j\le m$, where the $a_i$ are algebraically independent. For $0< s\le m+1$, write $h_s=H_{S_K/(IS_K,\ell_0,\ldots,\ell_{s-1})}(d)$ and $h_0=H_{\F_q[X_0,\ldots,X_m]/I}(d)$. Then $h_s\le g_d^{\circ s}(h_0)$.
\end{proposition}

\begin{proof}
Fix $1\le s\le m+1$ and set
\[
\mathfrak a=(I,\ell_1,\ldots,\ell_{s-1})\subseteq K[X].
\]
Set $K_0=\F_q(a_0^q,\ldots,a_m^q)\subseteq K$. The map
\[
 \Phi:K[X]\longrightarrow K_0[X],
 \qquad
 \sum_{\alpha}c_{\alpha}X^{\alpha}
 \longmapsto
 \sum_{\alpha}c_{\alpha}^qX^{\alpha},
\]
is a graded ring isomorphism. Since $I$ is defined over $\F_q$,
we have $\Phi(IK[X])=IK_0[X]$, while
$\Phi(\ell_j)=\ell_{j+1}$ for $0\le j<m$. Hence $\Phi$ induces an isomorphism
\[
 \frac{K[X]}{(I,\ell_0,\ldots,\ell_{s-2})}
 \;\cong\;
 \frac{K_0[X]}{(I,\ell_1,\ldots,\ell_{s-1})}.
\]
Extending scalars on the right
from $K_0$ to $K$ preserves these dimensions, and therefore
\[
 H_{K[X]/\mathfrak a}(d)=h_{s-1}.
\]

Use coordinates $Y_i=X_i$ for $0\le i<m$ and
$Y_m=\ell_0=\sum_{i=0}^m a_iX_i$. Continue to denote the transformed
ideal by $\mathfrak a$, and let
\(
B=\mathfrak a\big|_{Y_m=0}\subseteq K[Y_0,\ldots,Y_{m-1}].
\)
Then $H_{K[Y']/B}(d)=h_s$, where $Y'=(Y_0,\ldots,Y_{m-1})$. Indeed, this follows from the isomorphisms
\[
\frac{K[Y']}{B}
\cong
\frac{K[Y',Y_m]}{(\mathfrak a,Y_m)}
\cong
\frac{K[X]}{(I,\ell_0,\ldots,\ell_{s-1})}.
\]

We write $a=(a_0,\dots,a_m)$. Fix
$c=(c_0,\ldots,c_{m-1},0)\in K^{m+1}$ and set
\[
\lambda_c=Y_m-\sum_{i=0}^{m-1}c_iY_i.
\]
For nonzero $F\in\mathfrak a_d$, write
\begin{equation}\label{eq:exp_F_b_Y}
F(a,Y)=\sum_{r=j}^dY_m^rF_r(a,Y')\in K[Y],\qquad F_j\ne0.
\end{equation}
After multiplying $F$ by a nonzero polynomial in $a$ (clearing denominators), we may assume
that its expression in the original generators $I,\ell_1,\ldots,\ell_{s-1}$
has polynomial coefficients in $a$. The inverse coordinate change
introduces only powers of $a_m$ in the denominators; multiplying
by a further power of $a_m$, if necessary, ensures that all $F_r$
also have polynomial coefficients in $a$.



Write $I=(g_1,\ldots,g_r)$, where $g_\nu\in\F_q[X]$.
In the original $X$-coordinates, after clearing denominators
as above, we may write
\[
F(a,X)=\sum_{\nu=1}^r A_\nu(a,X)g_\nu(X)
+\sum_{u=1}^{s-1}C_u(a,X)\ell_u(a,X),
\]
where $A_\nu,C_u\in\F_q[a,X]$. Consider an indeterminate $t$ and substitute
$a_i\mapsto a_i+tc_i$ for $0\le i<m$, keeping $a_m$ fixed.
Since the $g_\nu$ do not
involve $a$, and
\(
\ell_u(a+tc,X)=\ell_u(a,X)+t^{q^u}\sum_{i=0}^{m-1}c_i^{q^u}X_i,
\)
substituting $a\mapsto a+tc$ gives
\[
\begin{aligned}
F(a+tc,X)
=\sum_{\nu=1}^r A_\nu(a+tc,X)g_\nu(X)
+\sum_{u=1}^{s-1}C_u(a+tc,X)\ell_u(a,X)
+\sum_{u=1}^{s-1}\sum_{i=0}^{m-1}
  C_u(a+tc,X)t^{q^u}c_i^{q^u}X_i.
\end{aligned}
\]
Every coefficient in $t$ (which is a polynomial in $X$) of the first two sums belongs to
$\mathfrak a$, since these sums are combinations of its
generators. Every power of $t$ occurring in the last sum
has exponent at least $q$.
Therefore, every coefficient of $t^j$ with $j<q$
in $F(a+tc,X)$ belongs to $\mathfrak a$.

Applying the change from $X$ to $Y$ with $a$ fixed leaves
$t$ unchanged, so it carries each coefficient into the same
ideal expressed in the $Y$-coordinates. In these coordinates,
the shifted polynomial is
\(
F\left(a+tc,Y',Y_m+t\sum_{i=0}^{m-1}c_iY_i\right),
\)
and its coefficients of $t^j$ with $j<q$ belong to $\mathfrak a$.
Setting $Y_m=0$, all these coefficients belong to $B$, since
this sends $\mathfrak a$ to $B$. From \eqref{eq:exp_F_b_Y},
the resulting polynomial is
\[
\sum_{r=j}^d t^r\left(\sum_{i=0}^{m-1}c_iY_i\right)^r
F_r(a+tc,Y').
\]
Its coefficient of $t^j$ is exactly
\begin{equation}\label{eq:lowestparameter}
\left(\sum_{i=0}^{m-1}c_iY_i\right)^jF_j(a,Y').
\end{equation}
Since $j\le d<q$, the polynomial in \eqref{eq:lowestparameter}
belongs to $B$.

Give $Y_i$ weight one for $0\le i<m$ and
$Y_m$ weight zero, taking the largest weight as the initial polynomial.
Weighted initial ideals preserve Hilbert functions. Indeed, let $\prec$ be a monomial order.
Denote by $\prec_w$ the monomial order that compares
$w$-weights first and breaks ties using $\prec$.
By \cite[Prop.~1.8]{Sturmfels1996},
\(
\ini_{\prec}\bigl(\ini_w(\mathfrak a)\bigr)
=
\ini_{\prec_w}(\mathfrak a).
\)
Since $\mathfrak a$ and $\ini_w(\mathfrak a)$ are homogeneous,
preservation of Hilbert functions under monomial initial
ideals gives
\[
\begin{aligned}
H_{K[Y]/\mathfrak a}(d)
=H_{K[Y]/\ini_{\prec_w}(\mathfrak a)}(d)
=H_{K[Y]/\ini_{\prec}(\ini_w(\mathfrak a))}(d)
=H_{K[Y]/\ini_w(\mathfrak a)}(d).
\end{aligned}
\]
Thus, we obtain
\begin{equation}\label{eq:initialdimension}
H_{K[Y]/\ini_w(\mathfrak a)}(d)
=H_{K[Y]/\mathfrak a}(d)=h_{s-1}.
\end{equation}

Then, looking at \eqref{eq:exp_F_b_Y}, for $F\in \mathfrak{a}_d$ we have
\(
\ini_w(F)=Y_m^jF_j.
\)
By \eqref{eq:lowestparameter}, substituting
$Y_m=\sum_{i=0}^{m-1}c_iY_i$ in this initial polynomial
gives an element of $B$.
The initial polynomials of the elements of $\mathfrak a_d$ span
$\ini_w(\mathfrak a)_d$. Indeed, multiplying a homogeneous
$F\in\mathfrak a$ by a monomial $M$ gives
$\ini_w(MF)=M\ini_w(F)$, so every contribution to degree $d$
of the initial ideal is a linear combination of such initial
polynomials. By the restriction calculation above, every
$G\in\ini_w(\mathfrak a)_d$ therefore satisfies
\(
G\left(Y',\sum_{i=0}^{m-1}c_iY_i\right)\in B.
\)
Since this substitution also sends $\lambda_c$ to zero,
it induces a surjective linear map
\[
\left(\frac{K[Y]}{(\ini_w(\mathfrak a),\lambda_c)}\right)_d
\rightarrow
\left(\frac{K[Y']}{B}\right)_d.
\]
Surjectivity follows because the substitution leaves every
polynomial in $Y'$ unchanged. Hence
\(
h_s\le H_{K[Y]/(\ini_w(\mathfrak a),\lambda_c)}(d)
\)
for every $c_0,\ldots,c_{m-1}\in K$.
Every linear form with nonzero $Y_m$ coefficient is a nonzero
scalar multiple of some $\lambda_c$. Such forms constitute
the nonempty Zariski-open subset
\[
\mathcal O=
\left\{
\sum_{i=0}^m\alpha_iY_i\in K[Y]_1:\alpha_m\ne0
\right\}
\]
of the space of linear forms. On the other hand,
Theorem~\ref{thm:green} gives a nonempty Zariski-open subset
$\mathcal V\subseteq K[Y]_1$ on which Green's bound holds
for $K[Y]/\ini_w(\mathfrak a)$. The intersection $\mathcal O\cap\mathcal V$
is nonempty, and since $K$ is infinite it contains a
$K$-rational linear form $\lambda$. Therefore, using
\eqref{eq:initialdimension},
we obtain
\[
\begin{aligned}
h_s
\le H_{K[Y]/(\ini_w(\mathfrak a),\lambda)}(d)
\le g_d\bigl(H_{K[Y]/\ini_w(\mathfrak a)}(d)\bigr)
=g_d(h_{s-1}).
\end{aligned}
\]
The proof is finished by iterating, using the monotonicity of $g_d$.
\end{proof}

\begin{corollary}\label{cor:flagvanishing}
Let $I\subseteq S=\F_q[X_0,\ldots,X_m]$ be a homogeneous ideal,
$1\le d<q$, and $1\le s\le m+1$. Let
$K=\F_q(a_0,\ldots,a_m)$, with algebraically independent $a_i$,
and consider $\ell_j=\sum_{i=0}^m a_i^{q^j}X_i$. If
\(
 H_{S/I}(d)<\binom{d+s}{d},
\)
then
\(
 \big(K[X_0,\ldots,X_m]/
       (IK[X_0,\ldots,X_m],\ell_0,\ldots,\ell_{s-1})\big)_d=0.
\)
\end{corollary}

\begin{proof}
Apply Proposition~\ref{prop:frobenius}
and~\eqref{eq:vanishthreshold}.
\end{proof}



\section{Footprint techniques}\label{sec:transfer}

For an ideal $J$ in a polynomial ring $R$ over a field and a monomial
order $\prec$, its \emph{footprint} is
\[
 \Delta_\prec(J)=\{M:M\text{ is a monomial and }
                         M\notin\ini_\prec(J)\}.
\]
The monomials in the footprint are called \emph{standard monomials}.
Their residue classes form a basis of $R/J$
\cite[Thm.~15.3]{eisenbud}.
In particular, when $R/J$ is zero-dimensional, its vector-space dimension is the cardinality
of the footprint. This falls within the context of the footprint bound, a technique that has been widely used to bound the number of common zeros of polynomials \cite{BDG19footprint,sanjose_conjecture}. To prove Theorem~\ref{thm:main}, we use additional information
beyond the initial ideal of $I(\PP^m(\F_q))$.
A suitable change of coordinates, using linear forms whose
coefficients are successive $q$-th powers, reveals additional
monomials in the initial ideal and strengthens the resulting
footprint bound; see \eqref{eq:fieldrelation}
and~\eqref{eq:initialcontainment}.

A subset $\Lambda\subseteq\N^m$ is called a \emph{decreasing set} if
$u\in\Lambda$ and $v\le u$ componentwise imply $v\in\Lambda$.
We write $\Lambda_{\le t}=\{u\in\Lambda:\sum_i u_i\le t\}$.
The next proposition gives us our main tool to bound the number of zeros. 


\begin{proposition}\label{prop:transfer}
Let $\emptyset \ne Z\subseteq\PP^m(\F_q)$, $0\le k\le m$,
and $1\le d<q$. Suppose that
\(
 H_Z(d)<\binom{d+k+1}{d}.
\)
Then there is a decreasing set $\Lambda\subseteq\{0,\ldots,q-1\}^m$
such that
\begin{equation}\label{eq:transfer}
 \card Z\le\card\Lambda+p_{k-1},
 \text{ and }
 \card{\Lambda_{\le t}}=H_Z(t),
\end{equation}
for $0\leq t <q$.
\end{proposition}

\begin{proof}
Let $K=\F_q(a_0,\ldots,a_m)$, with the $a_i$ algebraically independent,
and write $I_K=I(Z)K[X_0,\ldots,X_m]$. We use the linear forms
\[
 L_i=\sum_{u=0}^m a_u^{q^{k-i}}X_u,
\]
for $0\le i\le k$.
Since $a_0,\ldots,a_k$ are algebraically independent over $\F_q$,
they are linearly independent over $\F_q$. Hence the Moore
determinant criterion \cite[Lem.~3.51]{LidlNiederreiter1997} gives
\(
\det\bigl(a_u^{q^i}\bigr)_{0\le i,u\le k}\ne0,
\)
and these linear forms
are linearly independent over $K$.



For every $P\in Z$, we have $L_0(P)\ne0$, by the algebraic
independence of the $a_i$. Write $K[X]=K[X_0,\ldots,X_m]$.
Since
\(
I_K=\bigcap_{P\in Z}I(P)K[X],
\)
and each extended point ideal is prime (its quotient is still a domain) and does not contain $L_0$,
we obtain
\[
(I_K:L_0)
=\bigcap_{P\in Z}(I(P)K[X]:L_0)
=I_K.
\]
Set $A=K[X]/(I_K,L_0)$. The standard colon-ideal exact sequence is
\[
0\longrightarrow
\bigl(K[X]/(I_K:L_0)\bigr)(-1)
\xrightarrow{\ \cdot L_0\ }
K[X]/I_K
\longrightarrow A
\longrightarrow0.
\]
Using $(I_K:L_0)=I_K$ and the invariance of Hilbert functions
under extension of scalars, taking dimensions in each degree gives
\begin{equation}\label{eq:artinian}
H_A(t)=H_Z(t)-H_Z(t-1),
\end{equation}
where $H_Z(-1)=0$. Since $H_Z(t)=\card Z$ for all sufficiently
large $t$, the ideal $(I_K,L_0)$ is zero-dimensional, and we get
\[
\dim_K A
=\sum_{t\geq 0}H_A(t)
=\card Z.
\]

Complete $L_0,\ldots,L_k$ to a basis of $K[X_0,\dots,X_m]_1$ with $T_1,\ldots,T_c$, where $c=m-k$. Write $Y_i=L_i$, and let $J$ be the image of $I_K$
after this change of coordinates and setting $Y_0=0$.
Then we can write $A=R/J$ with
\(
 R=K[T_1,\ldots,T_c,Y_k,\ldots,Y_1].
\)
We consider graded reverse lexicographic order with
$T_1\succ \cdots\succ T_c\succ Y_k\succ \cdots\succ Y_1$.

Let $V=\sum_v b_vX_v$ be any linear form over $K$, and let
$\Phi(V)=\sum_v b_vX_v^q$. Since the points of $Z$ are $\F_q$-rational,
we have
\begin{equation}\label{eq:fieldrelation}
 \begin{split}
 VL_i^q-L_{i-1}\Phi(V)
 =\sum_{u,v}a_u^{q^{k-i+1}}b_v
           (X_u^qX_v-X_uX_v^q)
 \in I_K,\text{ for } 1\le i\le k.
 \end{split}
\end{equation}
Take $V$ to be one of $T_1,\ldots,T_c,Y_i,\ldots,Y_k$. We change variables in \eqref{eq:fieldrelation} to $T_1,\dots,T_c,Y_k,\dots,Y_1,Y_0$, and quotient by $Y_0$. The first term on the left-hand side of~\eqref{eq:fieldrelation} is the
monomial $VY_i^q$, and it contains none of $Y_1,\ldots,Y_{i-1}$.
Every monomial of the second term is divisible by $Y_{i-1}$.
Therefore, the first term is larger in graded reverse lexicographic
order, and the leading monomial is $VY_i^q$. Thus, 
$$
\sum_{i=1}^kY_i^q(T_1,\ldots,T_c,Y_i,\ldots,Y_k)
 \subseteq\ini_\prec(J).
$$

On the other hand, Corollary~\ref{cor:flagvanishing}, applied with
$s=k+1$, gives
\[
 (R/(J,Y_1,\ldots,Y_k))_d 
\cong
(K[X]/(I_K,L_0,L_1,\ldots,L_k))_d
=0.
\]
Consequently, for each monomial $M$ in the $T$ variables of degree $d$, we have $M\in (J,Y_1,\ldots,Y_k)_d$, that is, we can write $M=F+\sum_{i=1}^kY_i H_i$, where $F\in J_d$, $H_i\in R_{d-1}$. Thus, taking $G=-\sum_{i=1}^k Y_i H_i$ we have $M+G=F\in J$, and every term of $G$ involves a $Y$ variable. Therefore, $M$ is the leading monomial of $M+G$ and we have proved
\begin{equation}\label{eq:initialcontainment}
 (T_1,\ldots,T_c)^d+
 \sum_{i=1}^kY_i^q(T_1,\ldots,T_c,Y_i,\ldots,Y_k)
 \subseteq\ini_\prec(J).
\end{equation}

Let $\Delta=\Delta_\prec(J)$ be the footprint for the chosen order.
By \eqref{eq:initialcontainment}, every monomial in $\Delta$ has total degree less than $d$
in the variables $T_1,\ldots,T_c$.
In particular, the exponent for each variable $T$ is less than $q$. Suppose that a
monomial in $\Delta$ has some exponent for a $Y$ variable at least $q$, and let $i$
be the smallest index with this property. By~\eqref{eq:initialcontainment},
the exponent of $Y_i$ is exactly $q$, all the exponents for the $T$ variables are zero,
and the exponents of $Y_{i+1},\ldots,Y_k$ are zero. Hence every
standard monomial whose exponent vector lies outside
$\{0,\ldots,q-1\}^m$ belongs to
\[
 \mathcal E_{q,k}
 =\bigcup_{i=1}^k
 \{Y_1^{b_1}\cdots Y_{i-1}^{b_{i-1}}Y_i^q:
                     0\le b_j<q\}.
\]
This set has cardinality $1+q+\cdots+q^{k-1}=p_{k-1}$,
and all its monomials have degree at least $q$.

Identify monomials with their exponent vectors and set
$\Lambda=\Delta\cap\{0,\ldots,q-1\}^m$.
The standard monomials are closed under division, so $\Lambda$ is
a decreasing set.
Since $\card\Delta=\dim_K A=\card Z$ and
$\Delta\setminus\Lambda\subseteq\mathcal E_{q,k}$, we obtain
\[
\card Z
=\card\Lambda+\card{\Delta\setminus\Lambda}
\le\card\Lambda+\card{\mathcal E_{q,k}}
=\card\Lambda+p_{k-1},
\]
which gives the first inequality in~\eqref{eq:transfer}. Note that \(
H_A(j)=\dim_K(R/J)_j=|\Delta_j|,
\)
where $\Delta_j:=\{M\in\Delta:\deg M=j\}$. For $j<q$ every standard monomial in $\Delta_j$ has its exponent vector in $\Lambda$, i.e., 
\[
|\Lambda_{\le t}|
=\sum_{j=0}^{t}|\Delta_j|
=\sum_{j=0}^{t}H_A(j)
=H_Z(t),
\]
for $0\le t<q$, where the last equality follows from~\eqref{eq:artinian}.
\end{proof}

We can now prove the main result combining Proposition~\ref{prop:transfer} with Theorem~\ref{thm:affine}.

\begin{proof}[Proof of Theorem~\ref{thm:main}]
Let $Z$ be the common rational zero set of $r$ independent elements
of $S_d$, and let $h=N-r$. By~\eqref{eq:HZbound}, $H_Z(d)\le h$.
Let $k=\mu_d$. Since $h<N$, we have $0\le k\le m-1$,
and~\eqref{eq:macrange} gives
$H_Z(d)\le h<\binom{d+k+1}{d}$. By
Proposition~\ref{prop:transfer}, there exists a decreasing set
$\Lambda\subseteq\{0,\ldots,q-1\}^m$ such that
\[
 \card Z\le\card\Lambda+p_{k-1}
 \text{ and } \card{\Lambda_{\le d}}=H_Z(d)\le h.
\]
Fix an ordering $\gamma_0,\ldots,\gamma_{q-1}$ of $\F_q$, and set
\[
 Y_\Lambda=\{(\gamma_{u_1},\ldots,\gamma_{u_m}):u\in\Lambda\}.
\]
Then $\card{Y_\Lambda}=\card\Lambda$.
For $v\in\{0,\ldots,q-1\}^m$, consider the polynomial
\begin{equation}\label{eq:affineproducts}
 N_v(X_1,\ldots,X_m)
 =\prod_{i=1}^m\prod_{a=0}^{v_i-1}(X_i-\gamma_a).
\end{equation}
At the point indexed by $u$, the value of $N_v$ is nonzero if and
only if $v\le u$ componentwise. If $v\notin\Lambda$, then $N_v$
vanishes on $Y_\Lambda$: otherwise, some $u\in\Lambda$ would satisfy
$v\le u$, forcing $v\in\Lambda$ because $\Lambda$ is decreasing.

Take the polynomials $N_v$ with $v\notin\Lambda$ and
$\sum_i v_i\le d$. Their leading monomials
$X_1^{v_1}\cdots X_m^{v_m}$ are distinct, so they are linearly
independent. Since $d<q$, every nonnegative exponent vector of
total degree at most $d$ lies in $\{0,\ldots,q-1\}^m$.
Thus the number of these polynomials is
\[
 \binom{m+d}{d}-\card{\Lambda_{\le d}}\ge N-h=r.
\]
Applying Theorem~\ref{thm:affine} to any $r$ of them now yields
\[
 \card Z
 \le\card{Y_\Lambda}+p_{k-1}
 \le p_{k-1}+\sum_{i=1}^d\sq{\mu_i}
 =p_{\mu_d}+\sum_{i=1}^{d-1}\sq{\mu_i}.
\]
This proves the upper bound in~\eqref{eq:mainzeros}.
The reverse inequality was proved in~\cite[Thm. 2.3 and Cor. 2.6]{BDG22}.
It also follows independently from the point sets constructed in the
proof of Theorem~\ref{prop:generalbound} below. An explicit family
attaining the bound is given in Proposition~\ref{prop:newton}. Finally,
\eqref{eq:mainweights} follows from~\eqref{eq:weightszeros}.
\end{proof}

Wei duality also determines the weights in the following range of degrees. Together with Theorem~\ref{thm:main}, this result determines all the generalized
Hamming weights of $\PRM_q(d,2)$ for every positive degree $d$.

\begin{corollary}\label{cor:highdegree}
Let $m\ge1$, $(m-1)(q-1)<d<m(q-1)$, and let $e=m(q-1)-d$. The generalized Hamming weights of
$\PRM_q(d,m)$ are the integers in
\[
 \{1,\ldots,p_m\}\setminus
 \left\{p_m+1-d_s(\PRM_q(e,m)):
                  1\le s\le\binom{m+e}{e}\right\},
\]
where the weights of $\PRM_q(e,m)$ are
given by Theorem~\ref{thm:main}.
\end{corollary}
\begin{proof}
The result follows from S{\o}rensen's duality formula~\cite[Thm.~2]{Sorensen}, which gives
\(
 \PRM_q(d,m)^\perp=\PRM_q(e,m),
\)
and Theorem~\ref{thm:weiduality}.
\end{proof}
Note that the cases with $d\geq m(q-1)$ are trivial: for $d=m(q-1)$, the dual is spanned by the vector $(1,\dots,1)$, so the $r$th generalized Hamming weight is just $r+1$, and for $d>m(q-1)$ we have the full code. 

\section{A nested family attaining the weight hierarchy}\label{sec:chain}
Choosing the polynomials from the proof of \cite[Thm. 2.3]{BDG22} appropriately, we can find nested subspaces attaining the generalized Hamming weights, showing that $\PRM$ codes satisfy the chain condition, for $d<q$.
Let $1\le d<q$, and fix an ordering $\gamma_0,\ldots,\gamma_{q-1}$
of $\F_q$. For $0\le j\le m$, let $M=m-j$, and let
$\beta=(\beta_1,\ldots,\beta_{M+1})$ run through the nonnegative
tuples with $\sum_i\beta_i=d-1$, in descending lexicographic order.
Define
\begin{equation}\label{eq:newtonbasis}
 F_{j,\beta}
 =X_j^{\beta_{M+1}+1}
  \prod_{i=1}^{M}\prod_{a=0}^{\beta_i-1}
              (X_{j+i}-\gamma_aX_j).
\end{equation}
We order the polynomials first by increasing $j$ and, for each
fixed $j$, by descending lexicographic order of $\beta$.
Denote the resulting ordered family by $F_1,\ldots,F_N$.

\begin{proposition}\label{prop:newton}
Let $1\leq d <q$. The polynomials $F_1,\ldots,F_N$ form a basis of $S_d$. For every
$1\le r\le N$, the subcode
\[
 D_r=\Span_{\F_q}\{\ev(F_1),\ldots,\ev(F_r)\}
\]
has dimension $r$ and support of cardinality
$d_r(\PRM_q(d,m))$. Thus these codes satisfy the \emph{chain
condition}: the subcodes attaining the generalized Hamming weights
can be chosen nested, $D_1\subset\cdots\subset D_N$.
\end{proposition}

\begin{proof}
For $M\ge1$, for a fixed $j$, dehomogenizing the $F_{j,\beta}$ at $X_j=1$ gives
the Heijnen--Pellikaan family attaining the affine weight hierarchy
of order $d-1$ in $M$ variables, whose successive spans are nested
\cite[Thm.~5.10 and Cor.~5.12]{HP}. Each set $\{F_{j,\beta}\}_\beta$ is a basis of
$X_j\F_q[X_j,\ldots,X_m]_{d-1}$, and these spaces form a direct
sum equal to $S_d$. Evaluation is injective for $d<q$, so
$\dim D_r=r$. 
As in the proof of \cite[Thm.~2.3]{BDG22}, for every
$1\le r\le N$, the polynomials $F_1,\ldots,F_r$ have exactly
the number of $\F_q$-rational common zeros given by the
lower bound for $e_r(d,m)$ stated there. By
Theorem~\ref{thm:main}, this number is $e_r(d,m)$, proving
that every $D_r$ attains the corresponding generalized Hamming weight.
\end{proof}

This construction gives the following recursion when $d<q$.

\begin{corollary}\label{cor:recursion}
Let $1\le d<q$ and $m\ge1$. Then
\[
 d_r(\PRM_q(d,m))=
 \begin{cases}
 d_r(\RM_q(d-1,m)),&1\le r\le\binom{m+d-1}{d-1},\\[2mm]
 q^m+d_{r-\binom{m+d-1}{d-1}}(\PRM_q(d,m-1)),&\binom{m+d-1}{d-1}<r\le\binom{m+d}{d}.
 \end{cases}
\]
\end{corollary}

\begin{proof}
The set $\{F_{0,\beta}\}_\beta$ in~\eqref{eq:newtonbasis} has $\binom{m+d-1}{d-1}$ elements. Its
evaluations vanish at infinity, and on $X_0=1$ they form the
family attaining the affine weight hierarchy of order
$d-1$. This gives the first expression. Once that block is
complete, its support contains the whole affine chart, with $q^m$
points. The remaining blocks restrict at infinity to the same
construction in $\PP^{m-1}(\F_q)$, giving the second expression.
Equality follows from Proposition~\ref{prop:newton}.
\end{proof}

\section{An upper bound in every degree}\label{sec:general}

\subsection{Failure of the direct extensions}
The polynomial family in Section~\ref{sec:chain} extends naturally
to arbitrary degree by using the Heijnen--Pellikaan family for
$\RM_q(d-1,m-j)$ in each successive block, homogenized to degree
$d$ with respect to $X_j$. The resulting upper bound is obtained from the expression in
Corollary~\ref{cor:recursion} by replacing
$\binom{m+d-1}{d-1}$ with $\dim\RM_q(d-1,m)$ and
$\binom{m+d}{d}$ with $\dim\PRM_q(d,m)$. Equality can fail already for $d=q$.

\begin{example}\label{ex:generaldegree}
Take $q=3$, $m=2$, $d=3$, and $r=2$. Since $\dim\RM_3(2,2)=6$, the
direct extension would give $d_2(\PRM_3(3,2))=d_2(\RM_3(2,2))=5$, where
the last value follows from the Heijnen--Pellikaan formula~\cite{HP}.
However, the cubics $F_1=X_1^3-X_1X_0^2$ and $F_2=X_2^3-X_2X_0^2$ vanish
at all $9$ points of the chart $X_0=1$, and when $X_0=0$ their
evaluations are those of $X_1$ and $X_2$, since $a^3=a$ in $\F_3$. Hence
they span a two-dimensional subcode supported on the line $X_0=0$, of
cardinality $4$. One can check with \cite{sanjoseGHWsPackage,githubGHWs} that this subcode attains $d_2(\PRM_3(3,2))=4<5$.
\end{example}

The polynomials $X_i^q-X_iX_0^{q-1}$ explain the failure: they vanish
on the entire affine chart but can contribute nonzero evaluations
at infinity. Thus, an optimal subcode can be supported at infinity
before the first affine block has been filled. The following lemma
records this construction. 

\begin{lemma}\label{lem:lift}
Let $G=\lambda G'\in\F_q[X_0,\ldots,X_n]$ be homogeneous of degree
$e\ge1$, with $\lambda$ a linear form, and let
\[
 \tilde{G}=(\lambda^q-\lambda X_0^{q-1})\,G'
 =\prod_{\alpha\in\F_q}(\lambda-\alpha X_0)\,G'.
\]
Then $\tilde{G}$ is homogeneous of degree $e+q-1$, vanishes at every
point of $\PP^n(\F_q)$ with $X_0\ne0$, and $\tilde{G}(P)=G(P)$ for
every $P\in\PP^n(\F_q)$ with $X_0(P)=0$.
\end{lemma}

\begin{proof}
If $X_0(P)\ne0$, then $\lambda(P)/X_0(P)\in\F_q$, so one of the factors
$\lambda-\alpha X_0$ vanishes at $P$. If $X_0(P)=0$, then
$\tilde{G}(P)=\lambda(P)^qG'(P)=\lambda(P)G'(P)=G(P)$.
\end{proof}

Moreover, the chain condition can also fail in higher degree.

\begin{example}\label{ex:failurechain}
Take $C=\PRM_3(3,3)$, whose minimum distance is $9$.
For a projective plane $\mathcal H\subseteq\PP^3(\F_3)$, let $C_{\mathcal H}$
denote the subcode of $C$ supported on $\mathcal H$.
For $\mathcal H=\{X_0=0\}$, we have
$C_{\mathcal H}=\Span_{\F_3}\{\ev(X_i^3-X_iX_0^2):1\le i\le3\}$, since
the polynomials of degree at most $3$ vanishing on $\F_3^3$ are precisely
the linear combinations of the field equations $X_i^3-X_i$.
Thus $\dim C_{\mathcal H}=3$ and $\supp(C_{\mathcal H})=\mathcal H$, and
the same holds for every projective plane by a change of coordinates.
Since $\card{\mathcal H}=13$, we get $d_3(C)\leq 13$, and the Griesmer bound gives $d_3(C)=9+3+1=13$.

We claim that every three-dimensional subcode $D$ attaining $d_3(C)$
has a projective plane as its support. Every nonzero word of $D$ has
weight $9$, since
$\sum_{c\in D\setminus\{0\}}\operatorname{wt}(c)=2\cdot3^2\cdot13=26\cdot9$ and $D$ has 26 nonzero words, each of weight at least 9.
By the classification of minimum-weight words \cite[Thm.~4.3]{GL24},
each such word is supported on a projective plane with a line removed.
By inclusion-exclusion, any two of these nine-point supports intersect in at least $9+9-13=5$ 
points, whereas distinct projective planes meet in only four points.
Hence all these supports lie in the same plane $\mathcal H$, and
$D=C_{\mathcal H}$.

Now let $E\subseteq C$ be a four-dimensional subcode containing $D$.
Since $\dim C_{\mathcal H}=3$, a word of $E\setminus D$ restricts
nontrivially to the affine chart $\PP^3(\F_3)\setminus\mathcal H$. This
restriction belongs to $\RM_3(3,3)$, whose minimum distance is $6$, so
$\card{\supp(E)}\ge13+6=19$. On the other hand, the evaluations of the
four cubics $X_0X_1X_i$, $0\le i\le3$, are independent and have combined
support $\{X_0=1,\ X_1\ne0\}$, consisting of $18$ points.
Thus $d_4(C)\le18$, and no subcode attaining $d_3(C)$ can be
contained in a subcode attaining $d_4(C)$.
In particular, $C$ does not satisfy the chain condition.
\end{example}

This example shows that, in higher degree, the subcodes attaining
the generalized Hamming weights cannot always be obtained as
the successive spans of a single ordered family of polynomials.
This motivates the additional choices in the higher-degree construction
below, compared with the nested family in Section~\ref{sec:chain}
and~\cite[Conj.~2.4]{BDG22}.


\subsection{Filling affine spaces and their points at infinity}

The construction of Section~\ref{sec:chain} suggests arranging
common zeros by filling complete affine or projective spaces.
Recall that $H_Z(d)$ counts the independent conditions imposed by
vanishing on $Z$. If $k_d(m)=\dim\PRM_q(d,m)$, the subcode
vanishing on $Z\subseteq\PP^m(\F_q)$ has dimension
$k_d(m)-H_Z(d)$. Thus we seek large point sets that impose few
conditions on forms of degree $d$. 
For each $0\le n\le p_m$, we construct two sets of $n$ points
$Z_n^-$ and $Z_n^+$ in $\PP^m(\F_q)$ that lead to the
upper bound in Theorem~\ref{prop:generalbound}.


Fix the enumeration $\gamma_0,\ldots,\gamma_{q-1}$ of $\F_q$
used in Section~\ref{sec:chain}. List the points of
$\Aff^h(\F_q)$ by varying the last coordinate first, then the
next-to-last coordinate, and so on. For $0\le t\le q^h$, let
$L_h(t)$ be the set of the first $t$ points in this list.
For example, in $\Aff^2(\F_q)$ the list begins
\[
 (\gamma_0,\gamma_0),\,
 (\gamma_0,\gamma_1),\,\ldots,\,
 (\gamma_0,\gamma_{q-1}),\,
 (\gamma_1,\gamma_0),\,
 (\gamma_1,\gamma_1),\,\ldots.
\]

We now define two $n$-point sets $Z_n^-$ and $Z_n^+$
recursively. Put $Z_0^-=Z_0^+=\emptyset$ and, for $n>0$, write
\begin{equation}\label{eq:canonicaldigits}
 n=aq^s+b,\qquad 1\le a\le q-1,\qquad 0\le b<q^s.
\end{equation}
We have two cases:
\begin{enumerate}[label=\textup{(\roman*)}]
\item If $b\le p_{s-1}$, take $Y=L_m(aq^s)$.
      This is a union of $a$ complete parallel affine
      $s$-spaces, whose common intersection at infinity
      has $p_{s-1}$ points. Set
      \[
       Z_n^-=Y\cup Z_b^-,
       \qquad
       Z_n^+=Y\cup Z_b^+,
      \]
      where $Z_b^-$ and $Z_b^+$ are constructed recursively
      on this intersection.
\item If $b>p_{s-1}$, set
      \begin{equation}\label{eq:canonicalterminal}
       Z_n^-=L_{s+1}(n),\qquad
       Z_n^+=\mathcal H\cup L_{s+1}(n-p_s),
      \end{equation}
      where $\mathcal H$ is the hyperplane at infinity in
      $\PP^{s+1}(\F_q)$ and the affine sets lie in its complement.
\end{enumerate}
The construction involves at most $m$ recursive steps.
For $d\ge1$, let
\[
 h_d(n):=\min\{H_{Z_n^-}(d),H_{Z_n^+}(d)\}.
\]

\begin{example}\label{ex:planefilling}
Consider $\PP^2(\F_4)$. Its affine plane has $16$ points, and
its line at infinity has $5$ points. Each affine line has $4$
points. Parallel affine lines have a common point at infinity,
denoted by $Q$. 

For $n=9$, the sets $Z_9^-$ and $Z_9^+$ coincide: both consist
of two complete parallel affine lines together with $Q$,
or equivalently, two complete projective lines.
For $n=15$, the set $Z_{15}^-=L_2(15)$ consists of three
complete parallel affine lines and three points on the fourth.
The set $Z_{15}^+=\mathcal H\cup L_2(10)$ consists of the whole
line at infinity and ten affine points, arranged as two
complete parallel affine lines and two points on the next.

The smaller Hilbert function value can come from either set.
For $n=6$, the set $Z_6^-$ consists of one complete affine line
and two points on a parallel line, whereas $Z_6^+$ consists
of the whole line at infinity and one affine point.
Direct evaluation gives $H_{Z_6^-}(2)=5$, $H_{Z_6^+}(2)=4$,
whereas $H_{Z_{15}^-}(4)=13$, $H_{Z_{15}^+}(4)=14$.
\end{example}

\begin{lemma}\label{lem:lowerrealization}
Let $\Lambda\subseteq\{0,\ldots,q-1\}^m$ be a decreasing set, and set
\(
 Y_\Lambda=\{(\gamma_{u_1},\ldots,\gamma_{u_m}):u\in\Lambda\}.
\)
Then $\card{Y_\Lambda}=\card\Lambda$ and, for every integer $t\ge0$,
the space obtained by evaluating polynomials in
$\F_q[X_1,\ldots,X_m]$ of total degree at most $t$ on $Y_\Lambda$
has dimension $\card{\Lambda_{\le t}}$.
\end{lemma}

\begin{proof}
Reducing modulo $X_i^q-X_i$, for $1\le i\le m$, does not  change evaluations on $Y_\Lambda$.
We may therefore restrict to polynomials of total degree at most $t$
and degree less than $q$ in each variable.
The polynomials $N_v$ in~\eqref{eq:affineproducts}, with
$v\in\{0,\ldots,q-1\}^m$ and $\sum_i v_i\le t$, form a basis
of this space: their leading monomials are precisely the monomials
satisfying these degree bounds, and all other terms have smaller
total degree.
If $v\notin\Lambda$, then $N_v$ vanishes on $Y_\Lambda$, since
a nonzero value at the point indexed by $u\in\Lambda$ would imply
$v\le u$ and hence $v\in\Lambda$.
Thus the evaluations are spanned by those with $v\in\Lambda_{\le t}$.

For $v\in\Lambda_{\le t}$, evaluate $N_v$ on the points indexed by
$\Lambda_{\le t}$. Ordering the indices by any total order extending
the componentwise order gives a triangular matrix with nonzero
diagonal, since $N_v(P_v)\ne0$ for
$P_v=(\gamma_{v_1},\ldots,\gamma_{v_m})$.
These evaluation vectors are therefore linearly independent.
It follows that the dimension of the space of evaluations is exactly
$\card{\Lambda_{\le t}}$.
\end{proof}

Let $\sigma_q(i)$ denote the
sum of the base-$q$ digits of $i$, and set 
\[
 A(t,D)=\card{\{0\le i<t:\sigma_q(i)\le D\}}.
\]
The tuples $(u_1,\ldots,u_h)$ indexing the points
$(\gamma_{u_1},\ldots,\gamma_{u_h})$ of $L_h(t)$ form a
decreasing set, and in fact they are the base $q$ digit tuples of the integers $0,\dots,t-1$.
Thus Lemma~\ref{lem:lowerrealization} shows that $A(t,D)$ is the
dimension of the evaluations on $L_h(t)$ of polynomials in
$\F_q[X_1,\ldots,X_h]$ of total degree at most $D$, for
$0\le t\le q^h$.
In particular, $A(q^h,D)=\dim\RM_q(D,h)$ for $D\ge0$. Moreover, for $d\geq 1$, from \cite[Cor. 1]{SanJose} we have the following: 
\[
k_d(m)=A(q^m,d-1)+k_d(m-1) \text{ and } k_d(m)=\sum_{j=0}^m A(q^j,d-1).
\]
The next result, whose proof is given in
Appendix~\ref{app:hilbertfunctions}, gives a recursive formula for $h_d(n)$. For the recursion below, we consider $h_D(t)=0$ if $D\le0$.

\begin{proposition}\label{prop:canonicalfamilies}
For $d\ge1$ and $n>0$, write $n=aq^s+b$ as in
\eqref{eq:canonicaldigits}. Then
\begin{equation}\label{eq:canonicalprofile}
 h_d(n)=
 \begin{cases}
 A(aq^s,d)+h_{d-q+1}(b),&b\le p_{s-1},\\[1mm]
 \min\{A(n,d),\ k_d(s)+A(n-p_s,d-1)\},&b>p_{s-1}.
 \end{cases}
\end{equation}
\end{proposition}

We can now state the bound. 

\begin{theorem}\label{prop:generalbound}
For $1\le r\le k_d(m)$, let
\begin{equation}\label{eq:generalbound}
 \mathcal B_r(q,d,m):=p_m-
       \max\{0\le n\le p_m:h_d(n)\le k_d(m)-r\}.
\end{equation}
Then $d_r(\PRM_q(d,m))\le\mathcal B_r(q,d,m)$.
If $d<q$, equality holds.
\end{theorem}

\begin{proof}
Choose a maximizing $n$ and
$Z\in\{Z_n^-,Z_n^+\}$ with $H_Z(d)=h_d(n)$.
The subcode vanishing on $Z$ has dimension
$k_d(m)-h_d(n)\ge r$, and its support is contained in the
complement of $Z$. This gives the upper bound.

Now let $d<q$. For integers $\nu\ge0$, $1\le D<q$
and $q^\nu\le T<q^{\nu+1}$, we have the counting identity
\begin{equation}\label{eq:digitcount}
 A(T,D)=\binom{\nu+D}{D}+A(T-q^\nu,D-1).
\end{equation}
Indeed, the integers below $q^\nu$ whose digit sum is at most $D$
correspond to the nonnegative tuples of length $\nu$ whose sum
is at most $D$. Since $D<q$, each entry is automatically less
than $q$, so their number is $\binom{\nu+D}{D}$.
The remaining integers are $q^\nu+y$, with $0\le y<T-q^\nu$.
Since $y<(q-1)q^\nu$, adding $q^\nu$ introduces no carry, and
$\sigma_q(q^\nu+y)=1+\sigma_q(y)$. Thus, these integers contribute
$A(T-q^\nu,D-1)$ to $A(T,D)$.

Put $N=\binom{m+d}{d}=k_d(m)$. For $r=N$, both sides equal $p_m$, so assume $r<N$. Let $(\mu_1,\ldots,\mu_d)$ be the Macaulay tuple
of $N-r$, and let
\[
 \nu=\mu_d,\qquad T=\sum_{i=1}^d\sq{\mu_i},\qquad
 n=p_{\nu-1}+T.
\]
Since $0<N-r<N$, we have $0\le\nu\le m-1$.
Apply~\eqref{eq:digitcount} repeatedly to $A(T,d)$,
removing $q^{\mu_d},q^{\mu_{d-1}},\ldots$ in turn and lowering
the degree bound by one at each step.
At a step with $\mu_i\ge0$,
the remaining sum lies between $q^{\mu_i}$ and
$i q^{\mu_i}<q^{\mu_i+1}$, so the identity applies.
We stop when the remaining sum is zero. Thus
\[
 A(T,d)=\sum_{i=1}^d\binom{\mu_i+i}{i}=N-r.
\]

To apply~\eqref{eq:canonicalprofile}, observe that adding
$p_{\nu-1}$ to $T$ introduces no carries: each digit below
position $\nu$ receives at most $d-1$ contributions from $T$ (since $\nu=\mu_d$),
and one from $p_{\nu-1}$, giving at most $d<q$. We can write 
\[
n=
\underbrace{\#\{i:\mu_i=\nu\}}_{a}\,q^\nu
+
\underbrace{\left(
p_{\nu-1}+\sum_{\mu_i<\nu}\lfloor q^{\mu_i}\rfloor
\right)}_{b}.
\]
Consequently, the decomposition~\eqref{eq:canonicaldigits} of $n$
has $s=\nu$ and $b\ge p_{\nu-1}$.
In particular, $n<q^{\nu+1}\le q^m\le p_m$.

If $b=p_{\nu-1}$, then $T=aq^\nu$.
Since $d-q+1\le0$, the first case of~\eqref{eq:canonicalprofile}
gives $h_d(n)=A(T,d)$.
If $b>p_{\nu-1}$, the second case gives
\[
 \begin{aligned}
 h_d(n)
 \le k_d(\nu)+A(n-p_\nu,d-1)
 =\binom{\nu+d}{d}+A(T-q^\nu,d-1)
 =A(T,d).
 \end{aligned}
\]
Here $k_d(\nu)=\binom{\nu+d}{d}$ because $d<q$, and the last
equality follows from~\eqref{eq:digitcount}.

In both cases, one of the constructed sets $Z\subseteq\PP^m(\F_q)$
has $n$ points and $H_Z(d)=h_d(n)\le A(T,d)=N-r$.
The space of polynomials in $S_d$ vanishing on $Z$ therefore has
dimension at least $r$, proving $e_r(d,m)\ge n$.
Since $n=p_\nu+\sum_{i=1}^{d-1}\sq{\mu_i}$, the upper bound proved
in Theorem~\ref{thm:main} gives the reverse inequality, so
$e_r(d,m)=n$. Hence
\[
 \mathcal B_r(q,d,m)\le p_m-n=d_r(\PRM_q(d,m)),
\]
and equality follows from the upper bound $d_r(\PRM_q(d,m))\le\mathcal B_r(q,d,m)$.
\end{proof}

\begin{example}
   The bound recovers the value in Example~\ref{ex:generaldegree}. One can check that 
$h_3(9)=8=k_3(2)-2$, so
$\mathcal B_2(3,3,2)\le13-9=4$. 
\end{example}

By Theorem~\ref{prop:generalbound}, for $d<q$ the bound
$\mathcal B_r(q,d,m)$ coincides with the formula conjectured
by Beelen, Datta and Ghorpade~\cite[Conj.~2.4]{BDG22}
and proved in Theorem~\ref{thm:main}. We conjecture that the bound $\mathcal B_r(q,d,m)$
is sharp in every degree.

\begin{conjecture}\label{conj:general}
For every $m,d\ge1$ and $1\le r\le k_d(m)$,
\[
 d_r(\PRM_q(d,m))=\mathcal B_r(q,d,m).
\]
\end{conjecture}


\subsection{Some exact values}\label{subsec:lifts}
In this subsection, we provide some values of the weight hierarchy of projective Reed--Muller codes, in particular for degrees outside of the range of Theorem \ref{thm:main} and Corollary \ref{cor:highdegree}. 
First we show how to obtain subcodes of a projective Reed-Muller code supported on a smaller projective space using Lemma~\ref{lem:lift}.

\begin{lemma}\label{lem:generallifting}
Let $e\ge1$, $0\le t\le m-1$, and $h=m-t$. Then
\begin{equation}\label{eq:generallifting}
 d_r\bigl(\PRM_q(e+t(q-1),m)\bigr)
 \le d_r\bigl(\PRM_q(e,h)\bigr),
 \qquad 1\le r\le k_e(h).
\end{equation}
Moreover,
\[
 \mathcal B_r(q,e+t(q-1),m)=\mathcal B_r(q,e,m-t),
 \qquad 1\le r\le k_e(m-t).
\]
\end{lemma}

\begin{proof}
For a homogeneous polynomial $G\in\F_q[X_t,\ldots,X_m]$ of degree $e$,
apply Lemma~\ref{lem:lift} successively to each monomial and add the
resulting polynomials with the same coefficients as in $G$.
This gives a homogeneous polynomial of degree $e+t(q-1)$ whose
evaluations agree with those of $G$ on
$V=\{X_0=\cdots=X_{t-1}=0\}\cong\PP^h$ and are zero at every
rational point outside $V$. Thus every subcode of $\PRM_q(e,h)$
gives a subcode of $\PRM_q(e+t(q-1),m)$ with the same dimension
and support size. This proves~\eqref{eq:generallifting}.

For the identity for $\mathcal B_r$, it suffices to prove the case
$t=1$. Put $d=e+q-1$.
The recursion in~\eqref{eq:canonicalprofile} gives
\[
 h_d(q^m+b)=A(q^m,d)+h_e(b),
 \qquad 0\le b\le p_{m-1},
\]
and, taking $b=p_{m-1}$,
$k_d(m)=A(q^m,d)+k_e(m-1)$. Since $r\le k_e(m-1)$, we have
\[
 h_d(q^m)=A(q^m,d)\le k_d(m)-r.
\]
Thus the maximizing $n$ in~\eqref{eq:generalbound}
satisfies $n\ge q^m$. Writing $n=q^m+b$, we therefore have
\[
 h_d(n)\le k_d(m)-r
 \quad\Longleftrightarrow\quad
 h_e(b)\le k_e(m-1)-r.
\]
Together with $p_m-q^m=p_{m-1}$, this gives
$\mathcal B_r(q,e+q-1,m)=\mathcal B_r(q,e,m-1)$.
Iteration proves the claimed identity.
\end{proof}

\begin{remark}
Conjecture~\ref{conj:general} predicts that equality
always holds in~\eqref{eq:generallifting}.
It would be interesting to establish this equality
independently of the conjecture.
\end{remark}

For the cases arising from $e=1$ and $e=2$ below, S{\o}rensen's
minimum distance and the Griesmer bound give matching lower bounds.

\begin{proposition}\label{prop:lifts}
Let $0\le t\le m-1$ and $h=m-t$.
\begin{enumerate}[label=\textup{(\alph*)}]
\item For $1\le r\le h+1$,
 \[
  d_r\bigl(\PRM_q(1+t(q-1),m)\bigr)=q^h+q^{h-1}+\cdots+q^{h-r+1}.
 \]
\item If $q\ge3$, then
 \[
  d_r\bigl(\PRM_q(2+t(q-1),m)\bigr)=
  \begin{cases}
  q^h-q^{h-r},&1\le r\le h,\\
  q^h,&r=h+1.
  \end{cases}
 \]
\end{enumerate}
\end{proposition}

\begin{proof}
Let $V=\{X_0=\cdots=X_{t-1}=0\}\cong\PP^h$, with coordinates
$X_t,\ldots,X_m$. By Lemma~\ref{lem:generallifting}, it suffices
to construct subcodes of $\PRM_q(1,h)$
and $\PRM_q(2,h)$ with the required support sizes.

For the upper bound in (a), take $G=X_t,\ldots,X_{t+r-1}$. Their common
zeros in $V(\F_q)$ form a projective space of dimension $h-r$, so the support has $p_h-p_{h-r}=q^h+\cdots+q^{h-r+1}$
points. For (b), take $G=X_tX_{t+1},\ldots,X_tX_{t+r}$ if $r\le h$, and
$G=X_t^2,X_tX_{t+1},\ldots,X_tX_m$ if $r=h+1$. These vanish on
$V\cap\{X_t=0\}$, and on the chart $X_t=1$ of $V$, which is
$\Aff^h(\F_q)$ with coordinates $x_{t+1},\ldots,x_m$, they are the
functions $x_{t+1},\ldots,x_{t+r}$, respectively $1,x_{t+1},\ldots,x_m$.
Hence the supports have $q^h-q^{h-r}$ and $q^h$ points.

For the lower bounds, recall that if $1\le d\le m(q-1)$ and
$d-1=a(q-1)+b$ with $0\le b<q-1$, then
$d_1(\PRM_q(d,m))=(q-b)q^{m-a-1}$~\cite{Sorensen}. This gives $d_1=q^h$
in (a), and $d_1=(q-1)q^{h-1}$ in (b), where $q\ge3$ is used. Let $D$ be
an $r$-dimensional subcode. Deleting the coordinates outside $\supp(D)$
gives a code of length $\card{\supp(D)}$, dimension $r$ and minimum
distance at least $d_1$, so the Griesmer bound gives
\[
 \card{\supp(D)}\ge\sum_{i=0}^{r-1}\left\lceil\frac{d_1}{q^i}\right\rceil.
\]
In (a), this sum is $q^h+\cdots+q^{h-r+1}$ for $r\le h+1$. In (b), it is
$(q-1)(q^{h-1}+\cdots+q^{h-r})=q^h-q^{h-r}$ for $r\le h$, and
$q^h-1+1=q^h$ for $r=h+1$.
\end{proof}


\begin{remark}
For $t=1$, part (a) recovers the values of~\cite[Sec.~6]{BDG18} for $d=q$
and $r\le m$. For $q=3$, the degrees $1+2t$ and $2+2t$ cover every
$d\le 2m$, so the first $m-\lfloor (d-1)/2\rfloor+1$ weights of
$\PRM_3(d,m)$ are determined in every degree. Although $\mathcal B_r(q,d,m)$ is defined using Hilbert functions,
the explicit recursion in~\eqref{eq:canonicalprofile} makes it
straightforward to compute.
For all prime powers $q\le9$, $1\le m\le5$, and
$1\le d\le m(q-1)+1$, we checked that the bound agrees with
S{\o}rensen's minimum distance~\cite{Sorensen} and, in their
respective ranges, with the values of~\cite[Sec.~6]{BDG18},
Corollary~\ref{cor:highdegree}, and Proposition~\ref{prop:lifts}.
It also agrees with the complete weight hierarchies computed
using~\cite{sanjoseGHWsPackage} for $\PRM_q(d,2)$ with
$q\in\{3,4\}$ and $1\le d\le2(q-1)$. We also verified Conjecture~\ref{conj:general} for the complete
weight hierarchies of $\PRM_3(d,3)$ with $1\le d\le6$.
This follows by combining the preceding results with the bounds
in~\cite[Tables~3 and~6]{SanJose} and a computation
using~\cite{sanjoseGHWsPackage} confirming the remaining value
$d_7(\PRM_3(4,3))=16$.
\end{remark}






\section*{Declaration of AI}
ChatGPT-6 Astra was used to help organize the manuscript. The author laid out the proof strategy for the main result, and ChatGPT helped identify supporting references, in particular Green's hyperplane restriction theorem. ChatGPT has also been used to attempt to produce counterexamples for Conjecture \ref{conj:general}, to no avail.
All content, claims, and conclusions have been reviewed and remain the responsibility of the author.

\appendix
\section{Hilbert functions of the constructed point sets}\label{app:hilbertfunctions}
In this appendix, we provide the proof of Proposition~\ref{prop:canonicalfamilies}.

\begin{proof}[Proof of Proposition~\ref{prop:canonicalfamilies}]
First suppose $b\le p_{s-1}$, and consider $Y=\{\gamma_0,\ldots,\gamma_{a-1}\}\times\F_q^s$. Each affine space making up $Y$ is given by $X_1=\gamma_j$ in the chart $X_0=1$, so its projective closure satisfies
$X_1=\gamma_jX_0$, and setting $X_0=0$ forces $X_1=0$, i.e., their common intersection at infinity is
$L=\{X_0=X_1=0\}$. Let $W\subseteq L$ be either $Z_b^-$ or $Z_b^+$. We aim to compute $H_{Y\cup W}(d)$. Work in
$S=\F_q[X_0,\ldots,X_{s+1}]$ and consider the following exact sequence: 
\begin{equation}\label{eq:exact_sequence_recursive}
 0\longrightarrow \frac{I(Y)}{I(Y\cup W)}
 \longrightarrow \frac{S}{I(Y\cup W)}
 \longrightarrow \frac{S}{I(Y)}
 \longrightarrow 0,
\end{equation}
which gives
\[
 H_{Y\cup W}(d)
 =H_Y(d)+
 \dim_{\F_q}\left(\frac{I(Y)_d}{I(Y\cup W)_d}\right).
\]
Note that $\frac{I(Y)_d}{I(Y\cup W)_d}$ corresponds to the evaluations on $W$ of homogeneous polynomials of degree $d$ vanishing on $Y$.

Let $F$ be a homogeneous polynomial of degree $d$ vanishing
on $Y$. Divide $F(1,X_1,\ldots,X_{s+1})$ successively by
$\prod_{j=0}^{a-1}(X_1-\gamma_j)$ in $X_1$ and by
$X_i^q-X_i$ in each variable $X_i$, for $i\ge2$.
All these divisors vanish on $Y$, so the remainder does too.
Its degree is less than $a$ in $X_1$ and less than $q$ in
each other variable. Applying the univariate root bound
successively shows that the remainder is zero.
Homogenizing to degree $d$ gives
\[
 F=\left(\prod_{j=0}^{a-1}(X_1-\gamma_jX_0)\right)Q_1
   +\sum_{i=2}^{s+1}(X_i^q-X_iX_0^{q-1})Q_i,
\]
where $Q_1$ and the $Q_i$ for $i\ge2$ are homogeneous of
degrees $d-a$ and $d-q$, respectively.

On $L=\{X_0=X_1=0\}$, the first summand vanishes.
If $d<q$, the remaining summands are absent, so $F$ vanishes
on $L$. If $d\ge q$, then, for every $P\in L$,
\[
 F(P)=\sum_{i=2}^{s+1}X_i(P)^qQ_i(P)
     =\sum_{i=2}^{s+1}X_i(P)Q_i(P).
\]
Thus, the values of $F$ on $L$ are those of a homogeneous
polynomial of degree $e=d-q+1$ in $X_2,\ldots,X_{s+1}$. 

Conversely, every monomial $G$ of degree $e\ge1$ in these
variables is divisible by some $X_i$ with $i\ge2$. By
Lemma~\ref{lem:lift}, applied with $\lambda=X_i$, the polynomial
$\tilde{G}$ has degree $d$, vanishes on the entire affine chart, and
has the same values as $G$ on $L$.
Taking linear combinations therefore gives the evaluations
of every homogeneous polynomial of degree $e$ on $L$. Thus, for $e\ge1$, we have proved that the possible values on $L$ of polynomials of degree $d$ vanishing on $Y$ are exactly the values of the polynomials of degree $e$ on $L$, i.e., $\dim_{\F_q}\left(\frac{I(Y)_d}{I(Y\cup W)_d}\right)=H_W(e)$.
Since $H_Y(d)=A(aq^s,d)$, we obtain
\begin{equation}\label{eq:canonicalrank}
 H_{Y\cup W}(d)=A(aq^s,d)+
 \begin{cases}
 H_W(e),&e\ge1,\\
 0,&e\le0.
 \end{cases}
\end{equation}
This handles the case $b\leq p_{s-1}$ of the recursion \eqref{eq:canonicalprofile}.

Now suppose $b>p_{s-1}$, so the construction ends with the two
sets in~\eqref{eq:canonicalterminal}.
We compute their Hilbert functions separately, in $\PP^{s+1}(\F_q)$.

For the first choice, $Z_n^-=L_{s+1}(n)$, all $n$ points
are affine, so $H_{Z_n^-}(d)=A(n,d)$.
For the second choice, $Z_n^+=\mathcal H\cup L_{s+1}(v)$, where
\(
 v=n-p_s.
\)
Similar to the previous reasoning based on \eqref{eq:exact_sequence_recursive}, we first count the conditions on the whole hyperplane $\mathcal H$,
and then the additional conditions on the $v$ affine points.
The former contribute $k_d(s)$. To compute the latter, let
$F$ be homogeneous of degree $d$ and vanish on $\mathcal H$, and write
$F=X_0Q+R$, where $R$ does not involve $X_0$.
Since $R$ is homogeneous of positive degree and vanishes at
every projective point of $\mathcal H$, it vanishes at every vector
of $\F_q^{s+1}$.
Thus $F$ has the same affine values as $Q(1,X_1,\ldots,X_{s+1})$,
whose degree is at most $d-1$. Conversely, every affine
polynomial of degree at most $d-1$ is obtained by homogenizing
to degree $d-1$ and multiplying by $X_0$.
The same exact sequence, with $\mathcal H$ in place of $Y$, therefore gives
\[
 H_{Z_n^+}(d)=k_d(s)+A(v,d-1).
\]
Taking the minimum of the formulas for $Z_n^-$ and $Z_n^+$
gives the case $b>p_{s-1}$ of~\eqref{eq:canonicalprofile}.
Together with~\eqref{eq:canonicalrank}, this proves the
claimed Hilbert function formula by induction.
\end{proof}

\begin{thebibliography}{10}

\bibitem{munueraGHWhermitica}
A.~I. Barbero and C.~Munuera.
\newblock The weight hierarchy of {H}ermitian codes.
\newblock {\em SIAM J. Discrete Math.}, 13(1):79--104, 2000.

\bibitem{beelenGHWcartesian}
P.~Beelen and M.~Datta.
\newblock Generalized {H}amming weights of affine {C}artesian codes.
\newblock {\em Finite Fields Appl.}, 51:130--145, 2018.

\bibitem{BDG18}
P.~Beelen, M.~Datta, and S.~R. Ghorpade.
\newblock Maximum number of common zeros of homogeneous polynomials over finite fields.
\newblock {\em Proc. Amer. Math. Soc.}, 146:1451--1468, 2018.

\bibitem{BDG19footprint}
P.~Beelen, M.~Datta, and S.~R. Ghorpade.
\newblock Vanishing ideals of projective spaces over finite fields and a projective footprint bound.
\newblock {\em Acta Mathematica Sinica, English Series}, 35(1):47--63, 2019.

\bibitem{BDG22}
P.~Beelen, M.~Datta, and S.~R. Ghorpade.
\newblock A combinatorial approach to the number of solutions of systems of homogeneous polynomial equations over finite fields.
\newblock {\em Mosc. Math. J.}, 22(4):565--593, 2022.

\bibitem{Boguslavsky1997}
M.~Boguslavsky.
\newblock On the number of solutions of polynomial systems.
\newblock {\em Finite Fields Appl.}, 3(4):287--299, 1997.

\bibitem{sanjoseGHWNT}
E.~Camps-Moreno, H.~H. L\'opez, G.~L. Matthews, and R.~San-Jos\'e.
\newblock The weight hierarchy of decreasing norm-trace codes.
\newblock {\em Des. Codes Cryptogr.}, 93(7):2873--2894, 2025.

\bibitem{sanjose_complete_intersections}
E.~Camps~Moreno, F.~Salizzoni, and R.~San-José.
\newblock Generalized {H}amming weights of codes arising from complete intersections.
\newblock {\em ArXiv 2608.19978}, 2026.

\bibitem{sanjose_squarefree}
C.~Carvalho, H.~H. L\'opez, and R.~San-Jos\'e.
\newblock Cartesian square-free codes.
\newblock {\em Bol. Soc. Mat. Mex. (3)}, 32(3):Paper No. 125, 2026.

\bibitem{Caviglia}
G.~Caviglia.
\newblock A theorem of {Eakin} and {Sathaye} and {Green}'s hyperplane restriction theorem.
\newblock In {\em Commutative Algebra}, volume 244 of {\em Lect. Notes Pure Appl. Math.}, pages 1--5. Chapman \& Hall/CRC, Boca Raton, FL, 2006.

\bibitem{DG}
M.~Datta and S.~R. Ghorpade.
\newblock Number of solutions of systems of homogeneous polynomial equations over finite fields.
\newblock {\em Proc. Amer. Math. Soc.}, 145:525--541, 2017.

\bibitem{DGRemarks2017}
M.~Datta and S.~R. Ghorpade.
\newblock Remarks on the {Tsfasman--Boguslavsky} conjecture and higher weights of projective {Reed--Muller} codes.
\newblock In A.~Bassa, A.~Couvreur, and D.~Kohel, editors, {\em Arithmetic, Geometry, Cryptography and Coding Theory}, volume 686 of {\em Contemporary Mathematics}, pages 157--169. American Mathematical Society, Providence, RI, 2017.

\bibitem{eisenbud}
D.~Eisenbud.
\newblock {\em Commutative algebra}, volume 150 of {\em Graduate Texts in Mathematics}.
\newblock Springer-Verlag, New York, 1995.
\newblock With a view toward algebraic geometry.

\bibitem{EisenbudGreenHarris1993}
D.~Eisenbud, M.~Green, and J.~Harris.
\newblock Higher {Castelnuovo} theory.
\newblock In {\em Journ\'ees de g\'eom\'etrie alg\'ebrique d'Orsay -- Juillet 1992}, number 218 in Ast\'erisque, pages 187--202. Soci\'et\'e math\'ematique de France, 1993.

\bibitem{eisenbud_green_harris}
D.~Eisenbud, M.~Green, and J.~Harris.
\newblock Cayley-{B}acharach theorems and conjectures.
\newblock {\em Bull. Amer. Math. Soc. (N.S.)}, 33(3):295--324, 1996.

\bibitem{geilbezout}
O.~Geil and T.~H{\o}holdt.
\newblock Footprints or generalized {B}ezout's theorem.
\newblock {\em IEEE Trans. Inform. Theory}, 46(2):635--641, 2000.

\bibitem{GL24}
S.~R. Ghorpade and R.~Ludhani.
\newblock On the minimum distance, minimum weight codewords, and the dimension of projective {Reed--Muller} codes.
\newblock {\em Advances in Mathematics of Communications}, 18(2):360--382, 2024.

\bibitem{sanjoseSSCPRM}
P.~Gimenez, D.~Ruano, and R.~San-Jos\'{e}.
\newblock Subfield subcodes of projective {R}eed-{M}uller codes.
\newblock {\em Finite Fields Appl.}, 94:Paper No. 102353, 2024.
\newblock 46 pp.

\bibitem{Green}
M.~Green.
\newblock Restrictions of linear series to hyperplanes, and some results of {Macaulay} and {Gotzmann}.
\newblock In {\em Algebraic Curves and Projective Geometry (Trento, 1988)}, volume 1389 of {\em Lecture Notes in Mathematics}, pages 76--86. Springer, Berlin, 1989.

\bibitem{HP}
P.~Heijnen and R.~Pellikaan.
\newblock Generalized {Hamming} weights of {$q$}-ary {Reed--Muller} codes.
\newblock {\em IEEE Trans. Inform. Theory}, 44(1):181--196, 1998.

\bibitem{KaplanKim2025}
N.~Kaplan and J.-L. Kim.
\newblock Hulls of projective {Reed--Muller} codes.
\newblock {\em Des. Codes Cryptogr.}, 93(3):683--699, 2025.

\bibitem{Lachaud}
G.~Lachaud.
\newblock Projective {Reed--Muller} codes.
\newblock In {\em Coding Theory and Applications (Cachan, 1986)}, volume 311 of {\em Lecture Notes in Computer Science}, pages 125--129. Springer, Berlin, 1988.

\bibitem{Lachaud1990}
G.~Lachaud.
\newblock The parameters of projective {Reed--M{\"u}ller} codes.
\newblock {\em Discrete Math.}, 81(2):217--221, 1990.

\bibitem{LidlNiederreiter1997}
R.~Lidl and H.~Niederreiter.
\newblock {\em Finite Fields}, volume~20 of {\em Encyclopedia of Mathematics and its Applications}.
\newblock Cambridge University Press, Cambridge, second edition, 1997.

\bibitem{LinArtinian2026}
Y.~Lin.
\newblock {Artinian} reduction applied to the {Largest Intersection Problem}.
\newblock {\em ArXiv 2611.01273}, 2026.

\bibitem{LinDimension2026}
Y.~Lin.
\newblock A {$k$}-dimensional version of the largest intersection problem.
\newblock {\em ArXiv 2608.17771}, 2026.

\bibitem{LinSinghal2025}
Y.~Lin and D.~Singhal.
\newblock Largest zero-dimensional intersection of {$r$} degree {$d$} hypersurfaces.
\newblock {\em J. Pure Appl. Algebra}, 230(8):Article 108316, 2026.

\bibitem{sanjose_conjecture}
J.~Nardi and R.~San-Jos\'e.
\newblock Maximum number of zeroes of polynomials on weighted projective spaces over a finite field.
\newblock {\em Bull. Lond. Math. Soc.}, 58(6):Paper No. e70417, 2026.

\bibitem{sanjose_recursive_WPRM}
J.~Nardi and R.~San-José.
\newblock Structure of weighted projective {R}eed-{M}uller codes.
\newblock {\em ArXiv 2603.24397}, 2026.

\bibitem{sanjoseHullsPRM}
D.~Ruano and R.~San-Jos\'e.
\newblock Hulls of projective {R}eed-{M}uller codes over the projective plane.
\newblock {\em SIAM J. Appl. Algebra Geom.}, 8(4):846--876, 2024.

\bibitem{sanjoseHullvariationPRM}
D.~Ruano and R.~San-Jos\'e.
\newblock Quantum error-correcting codes from projective {R}eed–{M}uller codes and their hull variation problem.
\newblock {\em Journal of Algebra and Its Applications}, 24(13--14):Article 2541009, 2025.

\bibitem{SanJose}
R.~San-Jos{\'e}.
\newblock A recursive construction for projective {Reed--Muller} codes.
\newblock {\em IEEE Trans. Inform. Theory}, 70(12):8511--8523, 2024.

\bibitem{sanjoseGHWMPC}
R.~San-Jos\'e.
\newblock About the generalized {H}amming weights of matrix-product codes.
\newblock {\em Comput. Appl. Math.}, 44(4):Paper No. 186, 2025.

\bibitem{sanjoseGHWsPackage}
R.~San-Jos\'{e}.
\newblock An algorithm for computing generalized {H}amming weights and the {S}age package {GHWs}.
\newblock {\em ACM Trans. Math. Softw.}, 51(4):Article 26, 2025.
\newblock 20 pp.

\bibitem{githubGHWs}
R.~San-Jos\'{e}.
\newblock {GHWs}: A {S}age package for computing the generalized {H}amming weights of a linear code. {G}it{H}ub repository.
\newblock Available online: \url{https://github.com/RodrigoSanJose/GHWs}, 2025.

\bibitem{sanjose_decoding_PRM}
R.~San-Jos\'e.
\newblock Recursive decoding of projective {R}eed-{M}uller codes.
\newblock {\em IEEE Trans. Inform. Theory}, 72(5):2769--2779, 2026.

\bibitem{Serre}
J.-P. Serre.
\newblock Lettre {\`a} {M. Tsfasman}.
\newblock {\em Ast{\'e}risque}, 198--200:351--353, 1991.

\bibitem{SL}
D.~Singhal and Y.~Lin.
\newblock On a conjecture of {Beelen, Datta and Ghorpade} for the number of points of varieties over finite fields.
\newblock {\em Eur. J. Math.}, 11(3):Paper No. 58, 2025.

\bibitem{SongLuo2026}
Y.~Song and J.~Luo.
\newblock Hull parameters of projective {Reed--Muller} codes.
\newblock {\em Des. Codes Cryptogr.}, 94:Paper No. 57, 2026.

\bibitem{SongYue2026}
Y.~Song and Q.~Yue.
\newblock Recursive structure of hulls of {PRM} codes.
\newblock {\em ArXiv 2604.21808}, 2026.

\bibitem{Sorensen}
A.~B. S{\o}rensen.
\newblock Projective {Reed--Muller} codes.
\newblock {\em IEEE Trans. Inform. Theory}, 37(6):1567--1576, 1991.

\bibitem{Sturmfels1996}
B.~Sturmfels.
\newblock {\em {Gr{\"o}bner} Bases and Convex Polytopes}, volume~8 of {\em University Lecture Series}.
\newblock American Mathematical Society, Providence, RI, 1996.

\bibitem{vardyIntractability}
A.~Vardy.
\newblock The intractability of computing the minimum distance of a code.
\newblock {\em IEEE Trans. Inform. Theory}, 43(6):1757--1766, 1997.

\bibitem{Wei}
V.~K. Wei.
\newblock Generalized {Hamming} weights for linear codes.
\newblock {\em IEEE Trans. Inform. Theory}, 37(5):1412--1418, 1991.

\bibitem{Zanella}
C.~Zanella.
\newblock Linear sections of the finite {Veronese} varieties and authentication systems defined using geometry.
\newblock {\em Des. Codes Cryptogr.}, 13:199--212, 1998.

\end{thebibliography}


% \bibliographystyle{abbrv}
% \bibliography{references}
\end{document}
